\section{Related Work}\label{related_work}

Die Erkennung von Cyberangriffen in \iot{}-Netzwerken mittels ML ist ein aktives und
wachsendes Forschungsfeld.
Die folgenden Abschnitte gliedern die zehn verwendeten Quellen in vier thematische Cluster:
Datensaetze und Benchmarks, ML-basierte Ansaetze zur Angriffserkennung, Behandlung von
Klassenungleichgewicht sowie Erklaerbarkeit von \ids{}-Modellen.

\subsection{Datensaetze und Benchmarks fuer IoT-Sicherheit}

Die empirische Forschung im Bereich \iot{}-Sicherheit ist grundlegend auf realistische und
repraesentative Benchmarkdatensaetze angewiesen.
Neto et al.\ \cite{neto2023ciciot2023} haben mit dem CICIoT2023-Datensatz einen umfassenden
Benchmark entwickelt, der Netzwerkverkehr von 105 realen \iot{}-Geraeten unter 33
verschiedenen Angriffsbedingungen enthaelt.
Die Angriffe sind in sieben Kategorien organisiert: Distributed Denial of Service (DDoS),
Denial of Service (DoS), Mirai, Reconnaissance, Spoofing, Brute-Force und Web-based Attacks;
gemeinsam mit dem Benign-Verkehr ergeben sich die acht in dieser Arbeit klassifizierten
Kategorien.
Ein zentrales Merkmal des Datensatzes ist, dass die Angriffe ausschliesslich von \iot{}-Geraeten
gegen andere \iot{}-Geraete ausgefuehrt wurden, was eine besonders realitaetsnahe Abbildung
realer Bedrohungsszenarien ermoeglicht~\cite{neto2023ciciot2023}.
Der CICIoT2023-Datensatz weist ein ausgepraegtes Klassenungleichgewicht auf, das bei der
Modellbewertung methodische Herausforderungen mit sich bringt.

\subsection{ML-basierte Ansaetze zur Angriffserkennung in IoT-Netzwerken}

Hamedani et al.\ \cite{hamedani2025iotsurvey} liefern mit ihrer umfassenden
Uebersichtsarbeit einen systematischen Ueberblick ueber ML-basierte Sicherheitsloesungen fuer
\iot{}-Netzwerke aus den Jahren 2020 bis 2024.
Die Autoren klassifizieren ML-Techniken nach ihrer spezifischen Anwendung in der
\iot{}-Sicherheit und identifizieren offene Forschungsfragen, darunter insbesondere den Mangel
an interpretierbaren Modellen fuer ressourcenbeschraenkte Umgebungen~\cite{hamedani2025iotsurvey}.
Kikissagbe und Adda~\cite{kikissagbe2024review} ergaenzen diesen Ueberblick durch eine
systematische Analyse von supervised, unsupervised und hybriden ML-Ansaetzen fuer \ids{} in
\iot{}-Umgebungen.
Ihre Arbeit zeigt, dass baumbasierte Ensemble-Methoden konsistent zu den
leistungsstaerksten Ansaetzen gehoeren, gleichzeitig aber Fragen der Interpretierbarkeit
und Ressourceneffizienz weitgehend unbeantwortet bleiben~\cite{kikissagbe2024review}.

Auf dem CICIoT2023-Datensatz selbst haben Raturi et al.\ \cite{raturi2026exploratory} eine
vergleichende Evaluation mehrerer ML-Algorithmen durchgefuehrt.
Sie zeigen, dass baumbasierte Modelle und Ensemble-Methoden die hoechsten \ba{}-Werte
erreichen, wobei \gb{} in der Mehrklassenklassifikation den besten Wert unter den getesteten
Modellen erzielt~\cite{raturi2026exploratory}.
Die Autoren fuehren einen zweistufigen hierarchischen Klassifikationsansatz ein, der das starke
Klassenungleichgewicht des Datensatzes strukturell adressiert: In Stufe~1 wird zwischen
DoS/DDoS-Verkehr und Nicht-DoS-Verkehr unterschieden, in Stufe~2 erfolgt die
Mehrklassenklassifikation der verbleibenden Kategorien~\cite{raturi2026exploratory}.
Almahaqeri et al.\ \cite{almahaqeri2026gradient} evaluieren ebenfalls \gb{} auf dem
CICIoT2023-Datensatz und waehlen eine andere Strategie zur Behandlung des
Klassenungleichgewichts, naemlich stratifiziertes Undersampling.
Alharby~\cite{alharby2025ciciot} vergleicht auf demselben Datensatz mehrere ML-Modelle
einschliesslich \gb{} und Decision Trees unter Einsatz von \pca{} als
Dimensionsreduktionstechnik und dokumentiert zusaetzlich Trainings- und Vorhersagezeiten.

\subsection{Behandlung von Klassenungleichgewicht in IDS-Datensaetzen}

Das Klassenungleichgewicht stellt eine der zentralen methodischen Herausforderungen bei der
Entwicklung ML-basierter \ids{} dar.
Hosseini et al.\ \cite{hosseini2025imbalance} untersuchen systematisch die Auswirkungen
verschiedener Grade von Klassenungleichgewicht auf die Leistung von ML-Algorithmen in
\ids{}-Datensaetzen.
Ihre Ergebnisse zeigen, dass die einfache Accuracy als Bewertungsmetrik bei unausgewogenen
Datensaetzen zu systematisch verzerrten Urteilen fuehrt~\cite{hosseini2025imbalance}.
Raturi et al.\ \cite{raturi2026exploratory} greifen diese Problematik auf und schlagen die \ba{}
als fairere Bewertungsgrundlage vor.
Sie belegen empirisch am Beispiel des Hidden Markov Models, dass ein Modell trotz eines
Recall-Werts von 0{,}999 fuer DoS-Angriffe in der Mehrklassenklassifikation auf einen
Recall-Wert von 0{,}13 einbrechen kann~\cite{raturi2026exploratory}.
Almahaqeri et al.\ \cite{almahaqeri2026gradient} verfolgen einen alternativen Ansatz und setzen
stratifiziertes Undersampling ein, um das Klassenungleichgewicht zu reduzieren, anstatt eine
hierarchische Klassifikationsarchitektur zu verwenden.

\subsection{Erklaerbarkeit von ML-basierten IDS}

Der Einsatz von Methoden der erklaerbaren Kuenstlichen Intelligenz (XAI) in \ids{} ist ein
juengeres, aber rasch wachsendes Forschungsfeld.
Ogunseyi et al.\ \cite{ogunseyi2026xai_review} analysieren in einem PRISMA-basierten
systematischen Review 129 peer-reviewte Studien aus den Jahren 2018 bis 2025 und stellen fest,
dass \shap{} und Local Interpretable Model-agnostic Explanations (LIME) in 98~Prozent aller
untersuchten XAI-IDS-Studien eingesetzt werden~\cite{ogunseyi2026xai_review}.
Ein zentrales Ergebnis ihrer Analyse ist, dass die Integration von post-hoc-XAI-Methoden die
Erkennungsgenauigkeit der zugrundeliegenden Modelle nicht beeintraechtigt~\cite{ogunseyi2026xai_review}.
Mohale und Obagbuwa~\cite{mohale2025xai} evaluieren mehrere ML-Modelle, darunter Decision Trees
und XGBoost, gemeinsam mit \shap{}, LIME und ELI5 auf einem \ids{}-Datensatz und zeigen, dass
\shap{} die konsistentesten und zuverlaessigsten Erklaerungen liefert.
Hermosilla et al.\ \cite{hermosilla2025xai_forensic} vergleichen \shap{} und LIME auf
XGBoost- und TabNet-Modellen im Kontext der forensischen Intrusion Detection und fuehren drei
quantitative Metriken zur Bewertung der Erklaerungsqualitaet ein: Fidelitaet, Konsistenz und
Jaccard-Aehnlichkeit.
Ihre Ergebnisse zeigen, dass \shap{} eine hoehere Fidelitaet und Konsistenz erzielt als
LIME~\cite{hermosilla2025xai_forensic}.

\subsection{Identifizierte Forschungsluecken und Positionierung}

Die Analyse der bestehenden Literatur offenbart eine konsistente und direkt belegbare
Forschungsluecke: Keine der identifizierten Studien vereint gleichzeitig alle drei folgenden
Elemente auf dem CICIoT2023-Datensatz.
Erstens die Verwendung der \ba{} als primaere Bewertungsmetrik.
Zweitens die Integration von \shap{} als post-hoc-Erklaerungskomponente.
Drittens eine systematische Ablation Study, die den Einfluss der \pca{}-Dimensionsreduktion
auf Modellleistung und \shap{}-Feature-Importance quantifiziert.

Raturi et al.\ \cite{raturi2026exploratory} verwenden zwar die \ba{} und den
CICIoT2023-Datensatz, verzichten jedoch vollstaendig auf jede Form der
Erklaerbarkeitsanalyse.
Almahaqeri et al.\ \cite{almahaqeri2026gradient} und Alharby~\cite{alharby2025ciciot} arbeiten
ebenfalls auf demselben Datensatz, integrieren aber weder \ba{} als primaere Metrik noch eine
\shap{}-Komponente.
Mohale und Obagbuwa~\cite{mohale2025xai} sowie Hermosilla et al.\ \cite{hermosilla2025xai_forensic}
liefern fundierte \shap{}-Integrationen mit quantitativer Qualitaetsbewertung, arbeiten
jedoch nicht auf dem CICIoT2023-Datensatz.

Die vorliegende Arbeit schliesst diese Luecke, indem sie einen \gb{}-Klassifikator auf dem
CICIoT2023-Datensatz mit \ba{} als primaerer Bewertungsmetrik und einer vollstaendigen
\shap{}-Analyse auf beiden Klassifikationsstufen verbindet.
